% ==================================================================
% id:       seq-003
% topic:    סדרות
% subtopic: סדרות מונוטוניות וחסומות
% solution: full
% ==================================================================
\textbf{עבור כל אחת מהסדרות הבאות, קבעו האם היא:}
\begin{itemize}
    \item
    חסומה /
 מלעיל ולא מלרע /  מלרע ולא מלעיל  / לא  מלעיל ולא  מלרע.
\item
 מונוטונית עולה / עולה ממש /  יורדת / יורדת ממש/  אינה מונוטונית.
\end{itemize}
\bE\bMs{2}
\item
$1,2,2,3,3,3,4,4,4,4,\ldots$
\begin{solbox}
    הסדרה אינה חסומה מלעיל, אך כן חסומה מלרע למשל על ידי
    $1$.
    בנוסף הסדרה מונוטונית עולה, אך לא עולה ממש.
\end{solbox}
\item
$a_n=(-1)^n-n$
\hint{כתבו את האיברים הראשונים, והשוו בין איברים סמוכים לפי הזוגיות של $n$.}
\begin{solbox}
    ניתן לראות שהסדרה לא חסומה מלרע, אך כן חסומה מלעיל כי
    לכל
    $n\ge1$
    מתקיים
    $a_n=(-1)^n-n\le 1-n\le 1$.
    מבחינת מונוטוניות, הסדרה
    $-n$
    היא מונוטונית יורדת אבל הוספה של
    $(-1)^n$
    קצת הורס את המונוטוניות, כפי שניתן להיווכח מתוך האיברים הראשונים בסדרה שלנו:
    $$a_1=-1-1=-2,\, a_2=-2+1=-1,\, a_3=-3-1=-4,\, a_4=-4+1=-3,\ldots$$
    באופן פורמלי, ניתן להסתכל על ההפרש בין שני איברים סמוכים. מצד אחד, לכל $k\ge1$
    $$a_{2k}-a_{2k+1}=(1-2k)-(-1-(2k+1))=3>0,$$
    כלומר $a_{2k+1}<a_{2k}$, ולכן הסדרה אינה עולה.
    מצד שני, לכל $k\ge1$
    $$a_{2k}-a_{2k-1}=(1-2k)-(-1-(2k-1))=1>0,$$
    כלומר $a_{2k-1}<a_{2k}$, ולכן הסדרה גם אינה יורדת.
    לסיכום: הסדרה חסומה מלעיל ולא מלרע, ואינה מונוטונית.
\end{solbox}
\item
$b_n=\frac{1}{n^2+1}$
\begin{solbox}
    הסדרה חסומה מלעיל כי לכל
    $n\ge1$
    מתקיים
    $ n^2+1\ge 1$
    לכן
    $$b_n=\frac{1}{n^2+1}\le \frac{1}{1}=1$$
  ומהצד השני היא חסומה מלרע כי היא תמיד אי-שלילית.
  הסדרה מונוטונית יורדת ממש, כי לכל
  $n\ge1$
  מתקיים
  $$.n^2< (n+1)^2\Longrightarrow n^2+1< (n+1)^2+1\Longrightarrow b_n=\frac{1}{n^2+1}> \frac{1}{(n+1)^2+1}=b_{n+1}$$
\end{solbox}
\item
$c_n=\sqrt{n}-\lfloor\sqrt{n}\rfloor$
\hint{היעזרו ב-$0\le x-\lfloor x\rfloor<1$. למונוטוניות, בדקו את ערכי הסדרה כאשר $n$ הוא ריבוע שלם.}
\begin{solbox}
   ראשית, לכל מספר
   $x\in\rr$
   מתקיים
   $x\ge \lfloor x\rfloor$
   ולכן
   $c_n\ge 0$,
   אז היא חסומה מלרע.
   בנוסף, מתקיים תמיד
   $x\le \lfloor x\rfloor +1$
   ולכן
   $c_n\le 1$,
   אז הסדרה חסומה מלעיל ובסה״כ חסומה.
לגבי מונוטוניות, המצב יותר מסובך אז נחפש דברים פשוטים: נשים לב כי בכל אינדקס שהוא ״ריבוע שלם״, כלומר מהצורה
$n=k^2$,
מתקיים
$$c_n=\sqrt{k^2}-\lfloor \sqrt{k^2}\rfloor=k-\lfloor k\rfloor=k-k=0$$
אבל האיבר הבא אחריו בסדרה אינו
שווה
$0$,
 כי
 $\sqrt{n^2+1}$
 בוודאות אינו שלם אז מתקיים
 $\lfloor \sqrt{n^2+1}\rfloor< \sqrt{n^2+1}$,
 כלומר $c_n>0$.
 לסיכום, בסדרה יש אינסוף מופעים של
 $0$
 ומיד אחריהם מספר חיובי, לכן אינה מונוטונית כלל.

\end{solbox}
\eMs
\eE
